\section{集合的势}

\begin{definition}[基数，势，有限集，可数集，连续统 / 连续势]
    略
\end{definition}

\begin{theorem}
    $\forall E$, $2^{\overline{\overline{E}}} > \overline{\overline{E}}$
\end{theorem}
\begin{proof}
显然 $2^{\overline{\overline{E}}} \ge \overline{\overline{E}}$，下证两者不相等. 

假设 $2^{\overline{\overline{E}}} = \overline{\overline{E}}$，则 $\exists$ 双射 $T: E \to 2^E$. 构造集合 $A = \{ x \in E \mid x \notin T(x) \}$
\textcolor{orange}{(所有映射到不包含自己的子集的元素)}

$A \subseteq E \implies A \in 2^E$， 则 $\exists x_0$, s.t. $T(x_0) = A$. 

若 $x_0 \notin A \implies T(x_0) \implies x_0 \in A$ 矛盾！$\implies 2^{\overline{\overline{E}}} \neq \overline{\overline{E}} \implies 2^{\overline{\overline{E}}} > \overline{\overline{E}}$
\end{proof}

$\to$ 记 $2^{\overline{\overline{E}}}$ 为 $2^E$

\begin{corollary}
    $2^{\aleph_0} > \aleph_0$
\end{corollary}

\begin{theorem}
    可列个可列集的并是可列集.
\end{theorem}
$\aleph_0^2 = \aleph_0 \times \aleph_0 = \aleph_0 \xrightarrow{(n)} \aleph_0^n = \aleph_0 \quad (\forall n \ge 1)$

\begin{corollary}
    $\overline{\overline{\mathbb{Q}}} = \aleph_0$
\end{corollary}

\begin{theorem}
    $\aleph > \aleph_0$
\end{theorem}
\begin{proof}
$\overline{\overline{(0, 1)}} = \aleph$，只需证 $\overline{\overline{(0, 1)}} > \aleph_0$. \\
假设 $\overline{\overline{(0, 1)}} = \aleph_0$， 则可以把 $(0, 1)$ 所有数排成一列. \\
将其十进制小数表示为 $x_k = 0.a_{k1}a_{k2}\dots \quad (k=1, 2, \dots)$ \\
令 $y = 0.y_1y_2\dots$， 其中 $y_k = \begin{cases} 1, & a_{kk} \neq 1 \\ 2, & a_{kk} = 1 \end{cases}$ \\
\textcolor{orange}{$\to$ 让 $y$ 的第 $k$ 位与 $x_k$ 不同. 也可以 $y_k = 9 - a_{kk}$} \\
则 $\forall k \in \mathbb{N}^*$, $y \neq x_k$， 而 $y \in (0, 1)$， 矛盾！ \\
$\implies \overline{\overline{(0, 1)}} > \aleph_0$， 即 $\aleph > \aleph_0$.
\end{proof}

\begin{theorem}[Bernstein 定理]
    若 $\overline{\overline{A}} \le \overline{\overline{B}}$，$\overline{\overline{B}} \le \overline{\overline{A}}$，则 $\overline{\overline{A}} = \overline{\overline{B}}$
\end{theorem}

代数数，超越数 

{$\cdot$ 全体代数数构成一个域.} 

{$\cdot$ 一个数为代数数 $\iff$ 实部、虚部均为代数数.}

\begin{theorem}
    代数数集是可列集
\end{theorem}
\begin{proof}
$\forall n \ge 1$, $a_nx^n + a_{n-1}x^{n-1} + \dots + a_1x + a_0$ 的总数 $\le \aleph_0^{n+1} = \aleph_0$. \\
零点数 $\le n$ 对应的代数数集 $A_n$，$\overline{\overline{A_n}} = \aleph_0$. \\
代数数集 $A = \bigcup_{n=1}^\infty A_n$， 由定理 1.4.3， $\overline{\overline{A}} = \aleph_0$.
\end{proof}

\begin{theorem}
    \quad $\aleph = 2^{\aleph_0}$
\end{theorem}
\begin{proof}
(1) $\overline{\overline{(0, 1)}} \to 2^{\mathbb{N}^*}$, $x = \sum_{n=1}^\infty a_n 2^{-n} \quad (a_n \in \{0, 1\})$. $F(x) = \{ n \in \mathbb{N}^* \mid a_n = 1 \}$ 是从 $(0, 1)$ 到 $2^{\mathbb{N}^*}$ 的单射(若 $F(x_1) = F(x_2)$， 显然有 $x_1 = x_2$). 所以，$2^{\aleph_0} \ge \overline{\overline{(0, 1)}}$

(2) $2^{\mathbb{N}^*} \to (0, 1)$: $\quad E \subset \mathbb{N}^*$, $G(E) = 0.a_1a_2\dots$， 其中 $a_n = \begin{cases} 1, & n \in E \\ 0, & \text{其他} \end{cases}$ 是从 $2^{\mathbb{N}^*}$ 到 $(0, 1)$ 的单射$\implies \overline{\overline{(0, 1)}} \ge 2^{\aleph_0}$.由 Bernstein 定理，$2^{\aleph_0} = \overline{\overline{(0, 1)}}$， 即 $2^{\aleph_0} = \aleph$
\end{proof}

\vspace{1em}


\begin{notebox}
{\heiti 刘维尔数：} $x$：$\forall n \in \mathbb{N}^+$， $\exists p, q \in \mathbb{Z}$， $p > 1$，
s.t. $$0 < \left| x - \frac{q}{p} \right| < \frac{1}{p^n}$$
$\to$ 对于刘维尔数， $\forall n \in \mathbb{N}^+$，
满足上式的 $(p, q)$ 有无数多对. 
\end{notebox}


\vspace{1em}

\begin{notebox}
{\heiti 狄利克雷逼近定理：} $\forall x \in \mathsf{C}_{\mathbb{R}} \mathbb{Q}$，
有无数多对 $(p, q)$ 满足 $$0 < \left| x - \frac{q}{p} \right| < \frac{1}{p^2}$$

\end{notebox}

\vspace{1em}

\begin{notebox}
{\heiti 无理测度：}
$$\sup \left\{ \mu \mid \text{存在无数多对整数 } (p, q) \text{ 满足 } 0 < \left| x - \frac{q}{p} \right| < \frac{1}{p^\mu} \right\}$$
\end{notebox}

\vspace{1em}

\begin{theorem}[（关于有理逼近的）刘维尔定理]
    设 $n \ge 2$，无理数 $x$ 是某个 $n$ 次整系数多项式的零点. 则存在常数 $A > 0$， $\forall q \in \mathbb{Z}, p \in \mathbb{Z}^+$，有 $\left| x - \frac{q}{p} \right| > \frac{A}{p^n}$
\end{theorem}
\begin{proof}
设 $f(x) = a_n x^n + a_{n-1} x^{n-1} + \dots + a_1 x + a_0 = 0$，则 
$$f\left(\frac{q}{p}\right) = a_n \frac{q^n}{p^n} + a_{n-1} \frac{q^{n-1}}{p^{n-1}} + \dots + a_1 \frac{q}{p} + a_0
= \frac{1}{p^n} \left( a_n q^n + a_{n-1} q^{n-1} p + \dots + a_1 q p^{n-1} + a_0 p^n \right)$$

所以$\left| f\left(\frac{q}{p}\right) \right| \ge \frac{1}{p^n}$，即 
\begin{align*}
    \left| f\left(\frac{q}{p}\right) \right| &= \left| f\left(\frac{q}{p}\right) - f(x) \right|\\
    &= \left| a_n \left( \frac{q^n}{p^n} - x^n \right) + a_{n-1} \left( \frac{q^{n-1}}{p^{n-1}} - x^{n-1} \right) + \dots + a_1 \left( \frac{q}{p} - x \right) \right| \\
    &= \left| \frac{q}{p} - x \right| \cdot \left| \sum_{i=1}^n a_i \sum_{j=0}^{i-1} \frac{q^j}{p^j} x^{i-j-1} \right|
\end{align*}

考虑 $\left| \frac{q}{p} - x \right| < 1$ 的情形，此时
$$\left| \frac{q}{p} \right| \le \left| \frac{q}{p} - x \right| + |x| < 1 + |x|$$
\begin{align*}
    \implies \left| \sum_{j=0}^{i-1} \frac{q^j}{p^j} x^{i-j-1} \right| &\le \sum_{j=0}^{i-1} (1+|x|)^j |x|^{i-j-1} \overset{\lambda = 1+|x|}{=} \frac{\lambda^i - 1}{\lambda - 1} |x|^{i-1} \\
    \implies \left| \sum_{i=1}^n a_i \sum_{j=0}^{i-1} \frac{q^j}{p^j} x^{i-j-1} \right| &\le a \sum_{i=1}^n \frac{\lambda^i - 1}{\lambda - 1} |x|^{i-1} < \frac{a\lambda}{\lambda - 1} \sum_{i=1}^n (\lambda|x|)^{i-1} \\
    &= \frac{a\lambda}{\lambda - 1} \frac{(\lambda|x|)^n - \lambda|x|}{\lambda|x| - 1} \text{ 是一个常数，记作 } C，C > 0
\end{align*}
其中 $a = \max \{|a_1|, |a_2|, \dots, |a_n|\}$. 所以$ \left| \frac{q}{p} - x \right| \cdot C \ge \frac{1}{p^n}$ 

取 $A = \min \left\{ 1, \frac{1}{C} \right\}$，则 $\forall p \in \mathbb{Z}^+, q \in \mathbb{Z}$，即有 $\left| \frac{q}{p} - x \right| \ge \frac{A}{p^n}$
\end{proof}

\summary{本节可以大致分为两个话题，一是集合的势，二是代数数、超越数、有理逼近等问题. }

\summary{证明集合的势有关的命题的重要工具是Bernstein定理，通过在两个方向上分别构造单射来证明两个集合势相等.}